\chapter{Defining Your Mentor Identity\\นิยามตัวตนของ Mentor}

การ Mentoring ที่มีคุณภาพเริ่มก่อนการพบ Mentee กล่าวคือเริ่มจากการที่ท่านรู้ว่ากำลังเข้าไปในความสัมพันธ์ด้วยบทบาทใด อำนาจแบบใด และความตั้งใจแบบใด ความชัดเจนนี้สำคัญเป็นพิเศษในมหาวิทยาลัย ซึ่งบทบาทอาจารย์ ผู้ประเมิน ผู้บริหาร และผู้ทรงคุณวุฒิมักซ้อนทับกัน หากไม่แยกบทบาท Mentee อาจตอบสิ่งที่คิดว่าท่านอยากได้ยิน แทนที่จะสำรวจสิ่งที่ตนจำเป็นต้องเรียนรู้จริง ๆ

\learningobjectives{
  \item แยกบทบาท Mentor ออกจาก Teacher, Coach และ Consultant
  \item เลือกใช้หมวกทั้งสี่ของ Mentor ให้เหมาะกับสถานการณ์
  \item ระบุจุดแข็ง อคติ และขอบเขตของตนเองในฐานะ Mentor
}

\section{บทบาทที่ต่างกันด้วยเจ้าของคำตอบ}
ความแตกต่างสำคัญมิได้อยู่ที่ชื่อเรียก แต่อยู่ที่ว่าใครเป็นเจ้าของคำตอบและใครรับผิดชอบการตัดสินใจ Teacher ถ่ายทอดองค์ความรู้ที่มีโครงสร้าง; Coach ใช้คำถามเพื่อปลดล็อกผลงานหรือเป้าหมายเฉพาะ; Consultant วิเคราะห์และเสนอคำตอบจากความเชี่ยวชาญ; ส่วน Mentor เดินเคียงข้าง Mentee ในระยะยาว ใช้ทั้งคำถาม ประสบการณ์ และเครือข่าย โดยคงให้ Mentee เป็นเจ้าของทิศทาง

\begin{longtable}{P{.18\textwidth}P{.2\textwidth}P{.2\textwidth}P{.2\textwidth}}
\toprule
\textbf{บทบาท} & \textbf{จุดเน้น} & \textbf{เจ้าของคำตอบ} & \textbf{ลักษณะสัมพันธ์}\\
\midrule
Teacher & ความรู้ตามหลักสูตร & Teacher & มีโครงสร้างและการประเมิน\\
Coach & ศักยภาพต่อเป้าหมาย & ผู้รับการ Coach & เจาะจงและมุ่งผลลัพธ์\\
Consultant & การแก้ปัญหา & Consultant เสนอทางเลือก & ตามขอบเขตโครงการ\\
Mentor & การเติบโตของบุคคล & Mentee & ต่อเนื่องและอยู่บนความไว้วางใจ\\
\bottomrule
\end{longtable}

Mentor อาจสลับไปทำหน้าที่อื่นได้ แต่ควรบอกให้ชัด เช่น \enquote{ช่วงนี้ท่านต้องการให้ผมฟังในฐานะ Mentor หรือให้เสนอความเห็นเชิงเทคนิคในฐานะผู้เชี่ยวชาญ} ประโยคสั้น ๆ นี้คืนสิทธิในการเลือกให้ Mentee และลดความสับสนเชิงอำนาจ

\section{หมวกทั้งสี่ของ Mentor}
\subsection{Sponsor: เปิดประตู}
Sponsor ใช้ความน่าเชื่อถือของตนทำให้ Mentee ได้รับการมองเห็น เช่น เสนอชื่อเข้าร่วมโครงการ หรือแนะนำให้ผู้ตัดสินใจรู้จัก การสนับสนุนที่รับผิดชอบต้องแยกระหว่าง \enquote{การเปิดโอกาส} กับ \enquote{การรับประกันผล}

\subsection{Connector: เชื่อมเครือข่าย}
Connector ช่วยให้ Mentee เข้าถึงผู้เชี่ยวชาญ ผู้ใช้ หรือภาคอุตสาหกรรมที่เหมาะสม การเชื่อมโยงที่ดีระบุเหตุผลและความคาดหวังแก่ทั้งสองฝ่าย มิใช่เพียงส่งรายชื่อหรืออีเมลต่อกัน

\subsection{Challenger: ท้าทายสมมติฐาน}
Challenger ช่วยให้ Mentee มองเห็นสิ่งที่ยังไม่อยากเห็น โดยอาศัยข้อมูล คำถาม และความเคารพ ไม่ใช่การตัดสิน ในวัฒนธรรมที่ให้ความสำคัญกับความอาวุโส หมวกนี้ต้องใช้ด้วยความอ่อนโยน แต่ไม่ควรถูกละเลย

\subsection{Cheerleader: ประคองความเชื่อมั่น}
Cheerleader ไม่ได้บอกว่า \enquote{ทุกอย่างจะดีเอง} แต่ช่วยแยกความล้มเหลวของการทดลองออกจากคุณค่าของคน การเล่าความผิดพลาดของตนอย่างพอดีอาจทำให้ Mentee เห็นว่าความยากเป็นส่วนหนึ่งของการเติบโต

\begin{examplebox}
Mentee คนหนึ่งได้รับข้อเสนอแนะจากแหล่งทุนว่าโครงการยังไม่มีหลักฐานการใช้งานจริง Mentor อาจเริ่มด้วยหมวก Cheerleader เพื่อรับฟังความผิดหวัง จากนั้นใช้ Challenger ถามว่า \enquote{ข้อเสนอแนะใดมีข้อมูลรองรับมากที่สุด} ต่อด้วย Connector เพื่อเชื่อมผู้ใช้ทดลอง และใช้ Sponsor เมื่อโครงการพร้อมรับการมองเห็นรอบใหม่ การเปลี่ยนหมวกตามจังหวะสำคัญกว่าการมีหมวกที่ถนัดเพียงใบเดียว
\end{examplebox}

\section{ข้อตกลงเชิงบทบาท}
ก่อนเริ่มความสัมพันธ์ ท่านและ Mentee ควรตกลงเป้าหมาย ความถี่ ช่องทางติดต่อ การรักษาความลับ และขอบเขตของคำแนะนำ ควรชัดเจนด้วยว่า Mentor ไม่ใช่ผู้อนุมัติผลงาน และ Mentee สามารถไม่เห็นด้วยได้โดยไม่กระทบความสัมพันธ์

\begin{mentornote}
ก่อนให้คำแนะนำ ลองถามว่า \enquote{ท่านอยากให้ผมช่วยคิด ช่วยเชื่อมคน หรือเสนอประสบการณ์ตรง} คำถามนี้ช่วยให้ความช่วยเหลือตรงกับความต้องการจริง
\end{mentornote}

\begin{keytakeaway}
ตัวตนของ Mentor มิได้เกิดจากความอาวุโสเพียงอย่างเดียว แต่เกิดจากความสามารถในการใช้อำนาจ ประสบการณ์ และเครือข่ายเพื่อเพิ่มความสามารถในการตัดสินใจของ Mentee
\end{keytakeaway}

\begin{reflection}
ให้คะแนนตนเอง 1--5 ในหมวก Sponsor, Connector, Challenger และ Cheerleader จากนั้นเขียนหลักฐานหนึ่งชิ้นสำหรับคะแนนแต่ละข้อ เลือกหมวกที่ใช้น้อยที่สุด และกำหนดพฤติกรรมเล็ก ๆ ที่จะทดลองในการสนทนาครั้งถัดไป
\blankline\blankline\blankline
\end{reflection}

\chaptersummary{Mentor ต่างจากบทบาทสนับสนุนอื่นตรงที่รักษาเจ้าของคำตอบไว้กับ Mentee ท่านสามารถเพิ่มคุณภาพของความสัมพันธ์ได้ด้วยการประกาศบทบาทให้ชัด เลือกหมวกให้เหมาะกับจังหวะ และสร้างข้อตกลงที่คุ้มครองความไว้วางใจตั้งแต่ต้น}
